\abstract % Структурный элемент: РЕФЕРАТ

\keywords{НАГРУЗОЧНОЕ ТЕСТИРОВАНИЕ, РЕЛЯЦИОННЫЕ БАЗЫ ДАННЫХ, СУБД, ПРОИЗВОДИТЕЛЬНОСТЬ, СРАВНИТЕЛЬНЫЙ АНАЛИЗ, ВЕБ-ИНТЕРФЕЙС, МЕТРИКИ ПРОИЗВОДИТЕЛЬНОСТИ, ВИЗУАЛИЗАЦИЯ}

Объектом исследования являются реляционные системы управления базами данных и их поведение под контролируемой нагрузкой. Предметом исследования выступают методы, алгоритмы и инструментальные средства для автоматизированного нагрузочного тестирования, сбора метрик производительности и визуализации результатов с целью сравнительного анализа. Объектом разработки является программный комплекс для сравнительного нагрузочного тестирования реляционных СУБД. Цель работы~--- создание инструмента, обеспечивающего автоматизированное тестирование, сбор метрик производительности и их анализ через единый веб-интерфейс.

В работе применены методы системного и объектно-ориентированного проектирования, математической статистики для сопоставления прогонов и экспериментального исследования производительности СУБД. Разработана клиент-серверная система с веб-интерфейсом и поддержкой MySQL, PostgreSQL и MariaDB; реализованы формирование нагрузки по сценариям, отображение метрик в реальном времени, хранение истории прогонов, сравнительный анализ результатов, а также подготовка и восстановление состояния баз данных при тестах с операциями записи. Развёртывание системы выполняется средствами Docker Compose.

Работоспособность подтверждена экспериментально на двух наборах данных~--- Sakila и Brazilian E-Commerce~--- при сравнении трёх СУБД; проведено 12~прогонов с нулевой долей ошибок.

Практическая значимость работы заключается в том, что результаты тестирования можно использовать для сравнения СУБД, оценки ускорения после изменений и выявления снижения производительности.